Intelligence has often been discussed as an ability, a property, an accumulated body of knowledge, or a possession of correctness. Each of these interpretations treats intelligence as a stopping structure. This paper rejects that framing and formulates intelligence as a dynamical phenomenon that persists only under specific structural conditions.

The paper proposes three minimal axiomatic conditions: a connected state space, temporal ordering, and a global constraint $\Phi$. These conditions are not presented as a complete theory of intelligence, but as a lower bound on the kind of dynamics required for intelligence-like motion. They distinguish perception, recognition, knowledge, and intelligence as different structural layers. The paper then compares two radically different substrates, the reticulate dynamics of \emph{Physarum polycephalum} and the layer transitions of Transformer architectures, and reads both through the same structural conditions. In particular, the flow-reinforcement dynamics of the Physarum solver formulated by Tero, Kobayashi, and Nakagaki (2007) already exhibits a structure homologous to the emergence of $\Phi$ as a global constraint irreducible to local update rules.

From this structure, the Gettier problem is reinterpreted not as a defect in the definition of knowledge, but as a mismatch between a stopping structure and the dynamical property it is asked to guarantee. Accumulated correctness does not stabilize intelligence. Instead, it increases inferential connectivity and generates a branching explosion. Therefore intelligence cannot be stabilized within an individual system alone: no stopping operator exists internally.

The unavoidable boundary implied by this result is named the \emph{differential horizon of intelligence}. This is not a second or separate Differential Horizon beyond VED Vol.4. It is a local horizon-form: the same descriptive boundary appearing within the observational window of intelligence. The paper situates this horizon-form within a four-layer correspondence spanning the physical layer (VED Vol.4), the biological layer (Bio-IFGT), the consciousness layer, and the intelligence layer, while leaving the detailed structure of externalized stopping mechanisms to subsequent work.
